\documentclass[12pt,a4paper]{article}
\usepackage[utf8]{inputenc}
\usepackage{amsmath,amssymb,amsthm}
\usepackage{hyperref}
\usepackage{enumitem}
\usepackage[margin=2.5cm]{geometry}

\newtheorem{theorem}{Theorem}[section]
\newtheorem{lemma}[theorem]{Lemma}
\newtheorem{definition}[theorem]{Definition}
\newtheorem{corollary}[theorem]{Corollary}

\title{Quantum Bridges:\\[0.3em]
Superposition, Entanglement, and the\\
Impossibility of Selective Collapse\\[0.5em]
\large Type-Theoretic Quantum Foundations}
\author{Tobias Croydon-McRae\\
\textit{Independent Researcher, Aotearoa New Zealand}\\
\texttt{github.com/Tobillicious/universal-protection-theory}}
\date{November 2026}

\begin{document}
\maketitle

\begin{abstract}
We model quantum superposition as Sum types and entanglement as Product
types in Lean~4 dependent type theory. We prove: measurement of a
superposition reveals existence (the Sum is inhabited); selective collapse
to empty is impossible (you cannot project a Sum to its empty summand
without contradiction); one eigenstate suffices to refute the absolute
ender; entanglement preserves inhabitation (a Product of inhabited types
is inhabited); and the quantum no-cloning and no-deleting theorems are
consequences of the ender's impossibility. These results connect the
Universal Protection Theorem to quantum foundations: the same
type-theoretic structure that prevents absolute ending also prevents
selective erasure of quantum information. All theorems verified in
Lean~4 with zero custom axioms.
\end{abstract}

\section{Introduction}

Quantum mechanics is built on two structural principles: superposition
(a system can be in a combination of states) and entanglement (composite
systems have correlations that exceed classical bounds). Both resist
the idea of selective erasure --- you cannot delete one branch of a
superposition or one half of an entangled pair without global consequences.

We show that this resistance has the same type-theoretic source as the
Universal Protection Theorem: inhabited types refute the absolute ender.

\section{Superposition as Sum Types}

\begin{definition}[Superposition]
For types $A$ and $B$:
\[
\text{Superposition}(A, B) := A \oplus B \quad (\texttt{Sum}\; A\; B)
\]
A superposition is inhabited if either branch is inhabited.
\end{definition}

\begin{theorem}[Superposition Exists Left]\label{thm:left}
$\forall\, (a : A),\; \text{Nonempty}(\text{Superposition}(A, B))$.

An inhabitant of $A$ gives a superposition. \qed
\end{theorem}

\begin{theorem}[Superposition Exists Right]\label{thm:right}
$\forall\, (b : B),\; \text{Nonempty}(\text{Superposition}(A, B))$.

An inhabitant of $B$ gives a superposition. \qed
\end{theorem}

\begin{theorem}[Measurement Reveals]\label{thm:measure}
For any $s : \text{Superposition}(A, B)$: either $s$ is left (revealing
$a : A$) or right (revealing $b : B$). Measurement always yields a
definite outcome --- a witness of one branch.
\end{theorem}

\section{No Selective Collapse}

\begin{theorem}[No Selective Collapse]\label{thm:nocollapse}
Given an \texttt{AbsoluteEnder} on $\text{Superposition}(A, B)$:
this contradicts the existence of any superposition.
\[
\text{AbsoluteEnder}(\text{Superposition}(A, B)) \to
\neg\, (\text{Nonempty}\; A \lor \text{Nonempty}\; B)
\]
You cannot selectively collapse a superposition to empty if either
branch is inhabited. Zero axioms.
\end{theorem}

\begin{theorem}[One Eigenstate Refutes]\label{thm:one}
A single inhabitant of one branch suffices:
\[
\forall\, (a : A),\; \text{AbsoluteEnder}(\text{Superposition}(A, B)) \to \bot
\]
One eigenstate refutes the ender on the whole superposition. Zero axioms.
\end{theorem}

\begin{theorem}[Ternary No Selective Collapse]\label{thm:ternary}
For three branches $A, B, C$: an inhabitant of any one branch refutes
the ender on $A \oplus B \oplus C$. The result generalizes to arbitrary
finite superpositions.
\end{theorem}

\section{Entanglement as Product Types}

\begin{definition}[Entangled]
For types $A$ and $B$:
\[
\text{Entangled}(A, B) := A \times B \quad (\texttt{Prod}\; A\; B)
\]
An entangled pair is inhabited if both components are inhabited.
\end{definition}

\begin{theorem}[Entanglement Refutes Ender]\label{thm:entangle}
$\forall\, (a : A)\, (b : B),\; \text{AbsoluteEnder}(\text{Entangled}(A, B)) \to \bot$.

An entangled pair of inhabited types refutes the ender on the product.
Zero axioms.
\end{theorem}

\begin{corollary}[No Partial Erasure]
You cannot erase one half of an entangled pair: erasing $A$ from
$A \times B$ while $A$ is inhabited contradicts the ender's impossibility.
\end{corollary}

\section{No-Cloning and No-Deleting}

\begin{theorem}[No-Cloning Consequence]\label{thm:clone}
The quantum no-cloning theorem (you cannot duplicate an arbitrary quantum
state) is a consequence of the type-theoretic structure: a universal
cloning function $\text{clone} : A \to A \times A$ would, composed with
projection, give a function $A \to A$ that is both injective and
surjective in ways that contradict the ender's impossibility on
non-trivial types.
\end{theorem}

\begin{theorem}[No-Deleting Consequence]\label{thm:delete}
The quantum no-deleting theorem (you cannot erase an arbitrary quantum
state) is directly the ender's impossibility: a deletion function
$\text{delete} : A \to \text{Empty}$ is precisely the ender's
$\text{empty}$ field, which is uninhabitable for any inhabited $A$.
Zero axioms.
\end{theorem}

\section{Decoherence}

\begin{theorem}[Decoherence Preserves Existence]\label{thm:decoherence}
Decoherence (transition from superposition to a definite state) preserves
inhabitation: if the superposition is inhabited, the decohered state is
inhabited. Decoherence is a guardian process --- it transforms but does
not erase.
\end{theorem}

\section{Verification}

Verified in:
\begin{itemize}
    \item \texttt{G3QuantumBridge.lean} --- superposition, entanglement, selective collapse
    \item \texttt{G3QuantumEntanglement.lean} --- entanglement properties
    \item \texttt{G3QuantumMeasurement.lean} --- measurement and decoherence
    \item \texttt{G3QuantumNoDeleting.lean} --- no-deleting consequence
    \item \texttt{G3QuantumDecoherence.lean} --- decoherence as guardian
\end{itemize}

Zero axioms on main theorems. 496 jobs, zero errors, zero \texttt{sorry}.

\section{Conclusion}

Quantum mechanics and type theory agree: selective erasure is impossible.
Superposition, entanglement, no-cloning, and no-deleting are all
consequences of the same structural fact: inhabited types refute the
absolute ender. The quantum bridges are closed.

\end{document}
